\documentclass{beamer}

\usecolortheme{beaver}
\setbeamertemplate{navigation symbols}{}
\setbeamertemplate{itemize items}[square]
\setbeamercolor{block title}{bg=gray!20, fg=red!70!black}

\usepackage{booktabs}
\usepackage{fancyvrb}
\usepackage{hyperref}
\usepackage{tabularx}

\DefineVerbatimEnvironment{CodeBlock}{Verbatim}{
    fontsize=\tiny,
    frame=single,
    rulecolor=\color{gray},
    framesep=2mm
}

\title{INTERCAL in Racket}
\subtitle{Reader, Normalizer, and Macro-Compiled Runtime}
\author{Ethan Hawk \& Eva Augur}
\date{\today}

\begin{document}

\begin{frame}
  \titlepage
\end{frame}

\begin{frame}
  \frametitle{Project Goal}

  \begin{itemize}
    \item Implement a substantial, source-compatible subset of C-INTERCAL in Racket.
    \item Support real programs, not just toy syntax fragments.
    \item Compile INTERCAL into ordinary Racket by macro expansion instead of interpreting an external bytecode.
    \item Preserve difficult semantics: \texttt{COME FROM}, \texttt{NEXT}/\texttt{RESUME}, \texttt{ABSTAIN}, gerunds, and tape I/O.
  \end{itemize}

  \vspace{0.2cm}
  \begin{block}{Current Status}
    The implementation now runs multiple upstream-style programs including \texttt{pit/fib.i}, \texttt{pit/hanoi.i}, \texttt{pit/hello.i}, \texttt{pit/life.i}, \texttt{pit/beer.i}, \texttt{pit/rot13.i}, and \texttt{pit/unlambda.i} for its hello-world input.
  \end{block}
\end{frame}

\begin{frame}
  \frametitle{What the Project Implements}

  \scriptsize
  \begin{tabularx}{\textwidth}{@{}l>{\raggedright\arraybackslash}X@{}}
    \toprule
    Area & Implemented functionality \\
    \midrule
    Reader/frontend & Source cleaning, continuation-line merging, lexer, brag grammar, parse-tree normalization \\
    Core statements & Assignment, \texttt{NEXT}, \texttt{COME FROM}, \texttt{READ OUT}, \texttt{WRITE IN}, \texttt{STASH}, \texttt{RETRIEVE}, \texttt{IGNORE}, \texttt{REMEMBER}, \texttt{FORGET}, \texttt{RESUME}, \texttt{GIVE UP}, \texttt{TRY AGAIN}, \texttt{NOTHING} \\
    Control-state modifiers & \texttt{NOT}, \texttt{ONCE}, \texttt{AGAIN}, percentage/chance execution, politeness checking \\
    Dynamic control & Label abstention, gerund abstention/reinstatement, delayed \texttt{COME FROM}, counted abstain state \\
    Data model & Onespots, twospots, tails/hybrids with multidimensional arrays, packed \texttt{SUB}, \texttt{BY}-dimensioned arrays \\
    Operators & \texttt{MINGLE}, \texttt{SELECT}, unary AND/OR/XOR, width-aware rewrites, constant folding \\
    I/O & Numeric input/output plus C-INTERCAL-style array tape I/O \\
    Libraries & \texttt{syslib.i} integration and label-triggered library loading \\
    \bottomrule
  \end{tabularx}
\end{frame}

\begin{frame}
  \frametitle{Missing or Incomplete Features}

  \begin{itemize}
    \item \texttt{CREATE} is not implemented.
    \item Multithreading/backtracking features are not implemented.
    \item PIC-INTERCAL specific functionality is not implemented.
    \item TriINTERCAL and operand-overloading extensions are not implemented.
    \item OIL optimizer idiom ingestion is not implemented; only built-in conservative optimizations exist.
    \item Some upstream pit programs still diverge or misbehave, so compatibility is substantial but not complete.
  \end{itemize}

  \vspace{0.2cm}
  \begin{alertblock}{Interpretation of ``done''}
    The project is no longer just a parser. It is a working language implementation with documented semantic gaps.
  \end{alertblock}
\end{frame}

\begin{frame}[fragile]
  \frametitle{Pipeline Overview}

  \begin{enumerate}
    \item Clean source text and merge continuation lines.
    \item Tokenize and parse INTERCAL into a concrete syntax tree.
    \item Normalize that tree into a strict S-expression IR.
    \item Reader emits a Racket module that calls the backend macro layer.
    \item Macros expand the IR into explicit Racket state machines.
    \item Runtime helpers enforce INTERCAL semantics and produce output.
  \end{enumerate}

  \vspace{0.2cm}
  \textbf{Representative source program}

  \begin{CodeBlock}
(10) DO .1 <- #1
(20) PLEASE DO (40) NEXT
(30) DO GIVE UP
(40) DO RESUME #1
\end{CodeBlock}

\end{frame}

\begin{frame}[fragile]
  \frametitle{Reader Stage}

  \begin{columns}[T]
    \begin{column}{0.47\textwidth}
      \textbf{Reader responsibilities}
      \begin{itemize}
        \item Consume raw source text.
        \item Run shared cleaning logic.
        \item Parse and normalize.
        \item Produce a Racket \texttt{module} form.
      \end{itemize}
    \end{column}

    \begin{column}{0.50\textwidth}
      \begin{CodeBlock}
(define current-intercal-language-module-path
  (make-parameter #f))

(define intercal-module-source
  (variable-reference->module-source
   (#%variable-reference)))

(define (default-language-module-path src)
  (cond
    [(path? src)
     (path->string
      (find-relative-path
       (or (path-only src) (current-directory))
       intercal-module-source))]
    [else "intercal.rkt"]))

;; 2. THE READER FUNCTIONS
(define (intercal-read in)
  (syntax->datum (intercal-read-syntax #f in)))

(define (intercal-read-syntax src in)
\end{CodeBlock}

    \end{column}
  \end{columns}
\end{frame}

\begin{frame}[fragile]
  \frametitle{Lexer and Grammar}

  \begin{columns}[T]
    \begin{column}{0.48\textwidth}
      \begin{CodeBlock}
(define (tokenize in)
  (define str (port->string in))
  (define clean-str
    (string-replace
     (string-replace str "!" "'.")
     "DON'T" "DO NOT"))
  (define words
    (expand-packed-subscripts
     (regexp-match* TOKEN-RX clean-str)))
  (for/list ([w words])
    (cond
      [(regexp-match #px"^[0-9]+$" w) (token 'NUMBER (string->number w))]
      [(equal? w "<-") (token 'GETS w)]
      [(equal? w "SUB") (token 'SUB w)]
      [else (token (string->symbol w) w)])))
\end{CodeBlock}

    \end{column}
    \begin{column}{0.47\textwidth}
      \begin{itemize}
        \item Packed \texttt{SUB} forms are expanded before tokenization.
        \item The grammar keeps source structure explicit: labels, prefixes, postfix modifiers, and operator precedence.
        \item Arrays and dimensions are parsed uniformly so they can normalize into one IR.
      \end{itemize}
    \end{column}
  \end{columns}
\end{frame}

\begin{frame}[fragile]
  \frametitle{Normalization to S-Expressions}

  \begin{columns}[T]
    \begin{column}{0.48\textwidth}
      \textbf{Concrete parse tree}
      \begin{CodeBlock}
(program
 (line
  (label "(" 10 ")")
  (stmt
   (do-prefix "DO")
   (op
    (assign
     (var "." (ident 1))
     "<-"
     (expr (mingle (select (unary (postfix (primary "#" 1))))))))))
 (line
  (label "(" 20 ")")
  (stmt
   (do-prefix "PLEASE")
   (do-prefix "DO")
   (op (next (target "(" 40 ")") "NEXT"))))
 (line (label "(" 30 ")") (stmt (do-prefix "DO") (op (giveup "GIVE" "UP"))))
 (line
  (label "(" 40 ")")
  (stmt
   (do-prefix "DO")
   (op
    (resume
     "RESUME"
     (expr (mingle (select (unary (postfix (primary "#" 1)))))))))))

\end{CodeBlock}

    \end{column}
    \begin{column}{0.48\textwidth}
      \textbf{Normalized IR}
      \begin{CodeBlock}
(sick-program/syslib
 (10 (do (assign |.1| (mesh 1))))
 (20 (please (next 40)))
 (30 (do (give-up)))
 (40 (do (resume (mesh 1)))))

\end{CodeBlock}

    \end{column}
  \end{columns}
\end{frame}

\begin{frame}[fragile]
  \frametitle{Normalizer Implementation}

  \begin{columns}[T]
    \begin{column}{0.60\textwidth}
      \input{generated/normalizer-snippet.tex}
    \end{column}
    \begin{column}{0.35\textwidth}
      \begin{itemize}
        \item Collapse parse scaffolding.
        \item Rewrite operators into stable symbolic forms.
        \item Normalize modifiers and control targets.
        \item Emit one IR understood by the backend.
      \end{itemize}
    \end{column}
  \end{columns}
\end{frame}

\begin{frame}[fragile]
  \frametitle{Reader Output Module}

  \begin{itemize}
    \item After normalization, the reader does not interpret the program.
    \item It emits a normal Racket module that exports \texttt{intercal-main}.
    \item That module body calls the macro-based backend built around \texttt{sick-program}.
  \end{itemize}

  \begin{CodeBlock}
(module intercal-mod "intercal.rkt"
  (provide intercal-main)
  (define (intercal-main)
    (call-with-values
     (lambda ()
       (sick-program/syslib
        (10 (do (assign |.1| (mesh 1))))
        (20 (please (next 40)))
        (30 (do (give-up)))
        (40 (do (resume (mesh 1))))))
     (lambda ignored (void))))
  (module+ main (intercal-main)))

\end{CodeBlock}

\end{frame}

\begin{frame}[fragile]
  \frametitle{\texttt{sick-program-core}: Macro Compilation}
  \scriptsize

  \begin{columns}[T]
    \begin{column}{0.57\textwidth}
      \begin{CodeBlock}
(define-syntax (sick-program-core stx)
  (syntax-parse stx
    [(_ (ln:integer lbl mod pct:integer
         is-not:boolean is-once:boolean is-again:boolean op) ...)
     (define ops (syntax->list #'(op ...)))

     (define all-vars
       (remove-duplicates
        (filter intercal-var?
                (flatten (map syntax->datum ops)))))

     (define grouped-come-froms
       (let ([h (make-hash)])
         (for-each
          (lambda (l-ln l-op)
            (syntax-parse l-op
              [((~datum come-from) target)
               (define t (eval-label-target #'target))
               (hash-set! h t (cons (syntax-e l-ln)
                                    (hash-ref h t '())))]
              [_ (void)]))
          (syntax->list #'(ln ...)) ops)
         (hash-map h cons)))

     (define gerund->lns-map
       (build-gerund-map #'(ln ...) ops))
     ...]))
\end{CodeBlock}

    \end{column}
    \begin{column}{0.38\textwidth}
      \begin{itemize}
        \item Collect variables from normalized ops.
        \item Precompute label maps and gerund maps.
        \item Generate one dispatch clause per runtime line.
        \item Emit explicit runtime tables for abstain, ignore, and control flow.
      \end{itemize}
    \end{column}
  \end{columns}
\end{frame}

\begin{frame}[fragile]
  \frametitle{Expanded Runtime Shape}

  \begin{itemize}
    \item Expansion makes state explicit: variables, stash stacks, NEXT stack, debug state, and output accumulator all become ordinary Racket bindings.
    \item This is the project’s compilation technique: compile INTERCAL into a Racket state machine rather than interpret strings at runtime.
  \end{itemize}

  \input{generated/sample-expanded.tex}
\end{frame}

\begin{frame}[fragile]
  \frametitle{Control-Flow Semantics}

  \begin{itemize}
    \item \texttt{NEXT} pushes a return target; \texttt{RESUME} removes one or more entries and jumps to the last removed target.
    \item \texttt{COME FROM} is resolved after the triggering line executes, with delayed behavior for resumed \texttt{NEXT} entries.
    \item \texttt{ABSTAIN} and \texttt{REINSTATE} are modeled with counted runtime state so overlapping abstentions compose correctly.
    \item \texttt{ONCE} and \texttt{AGAIN} mutate line state after execution or skip.
  \end{itemize}

  \vspace{0.2cm}
  \begin{block}{Why this matters}
    The hard part of INTERCAL is not parsing the keywords. It is compiling bizarre, self-modifying control flow into something predictable enough for Racket to run.
  \end{block}
\end{frame}

\begin{frame}[fragile]
  \frametitle{Conservative Optimizations}
  \scriptsize

  \begin{columns}[T]
    \begin{column}{0.58\textwidth}
      \input{generated/optimizer-snippet.tex}
    \end{column}
    \begin{column}{0.37\textwidth}
      \begin{itemize}
        \item Remove abstain checks for lines that cannot be abstained.
        \item Remove ignore-table lookups for variables that are never ignored.
        \item Remove \texttt{COME FROM} dispatch lookups for labels that cannot be hijacked.
        \item Use arithmetic/bitwise implementations for hot operators instead of list-based bit manipulation.
      \end{itemize}
    \end{column}
  \end{columns}
\end{frame}

\begin{frame}[fragile]
  \frametitle{Evidence: Real Programs}

  \begin{columns}[T]
    \begin{column}{0.52\textwidth}
      \textbf{Programs currently exercised in the repository}
      \begin{itemize}
        \item \texttt{pit/fib.i}
        \item \texttt{pit/hanoi.i}
        \item \texttt{pit/onceagain.i}
        \item \texttt{pit/hello.i}
        \item \texttt{pit/ignorret.i}
        \item \texttt{pit/beer.i}
        \item \texttt{pit/life.i}
        \item \texttt{pit/rot13.i}
        \item \texttt{pit/unlambda.i} (hello-world input)
      \end{itemize}
    \end{column}
    \begin{column}{0.44\textwidth}
      \textbf{Example program}
      \input{generated/hello-snippet.tex}
    \end{column}
  \end{columns}
\end{frame}

\begin{frame}
  \frametitle{User-Facing Documentation}

  \begin{itemize}
    \item We added Scribble documentation for language usage, supported features, semantic gaps, and the compilation pipeline.
    \item The documentation explains both current local use:
      \texttt{\#lang reader "intercal.rkt"}
    \item and the package-oriented goal:
      \texttt{\#lang intercal}
    \item The same documentation also serves as the packaging checklist for publishing to the Racket package catalog.
  \end{itemize}
\end{frame}

\begin{frame}
  \frametitle{What Is Needed for \texttt{\#lang intercal}?}

  \begin{enumerate}
    \item Expose a real \texttt{intercal/lang/reader.rkt} collection path.
    \item Make the package a multi-collection package if it provides both repository-root modules and an \texttt{intercal} collection.
    \item Add Scribble metadata to \texttt{info.rkt} so the package builds docs cleanly.
    \item Publish the repository to the package catalog with a stable source URL.
    \item Keep tests for both local \texttt{\#lang reader} usage and installed \texttt{\#lang intercal} usage.
  \end{enumerate}
\end{frame}

\begin{frame}
  \frametitle{Takeaways}

  \begin{itemize}
    \item The project implements an honest compiler pipeline: source text to CST to normalized IR to macro-expanded Racket.
    \item The main technical contribution is semantic encoding of INTERCAL’s dynamic control flow in ordinary Racket.
    \item The implementation is now capable of running nontrivial upstream programs, including a language implementation built on top of another esolang implementation.
    \item The next work is completeness: remaining extensions, more pit coverage, and packaging polish.
  \end{itemize}
\end{frame}

\end{document}
